\section*{अध्याय \ref{ch:b1:rate-laws} --- रासायनिक बलगतिकी: दर नियम}

\begin{solution}{exo:b1:rate-laws:1}
$v = \frac14 \times 2.0 \times 10^{-4} = \qty{5.0e-5}{mol.L^{-1}.s^{-1}}$;
डाइऑक्सीजन उसी दर, \qty{5.0e-5}{}, से बनती है; \ce{N2O5}
$2v = \qty{1.0e-4}{mol.L^{-1}.s^{-1}}$ की दर से लुप्त होता है।
\end{solution}

\begin{solution}{exo:b1:rate-laws:2}
कोटि 0: \unit{mol.L^{-1}.s^{-1}}; कोटि 1: \unit{s^{-1}}; कोटि 2:
\unit{L.mol^{-1}.s^{-1}}; कोटि 3: \unit{L^2.mol^{-2}.s^{-1}}।
\end{solution}

\begin{solution}{exo:b1:rate-laws:3}
$t_{1/2} = \ln2/k = 0.693/(3.0 \times 10^{-3}) = \qty{231}{s}$।
\qty{600}{s} के बाद $\eu^{-kt} = \eu^{-1.8} = 0.165$: 16.5\,\% शेष।
\end{solution}

\begin{solution}{exo:b1:rate-laws:4}
$1/a_0$ के समानुपाती \omterm{def:b1:rate-laws:half-life}{अर्ध-आयु} का अर्थ है कोटि 2। यदि $a_0$ को आधा करने से
\omterm{def:b1:rate-laws:half-life}{अर्ध-आयु} आधी हो जाए, तो $t_{1/2} \propto a_0$: कोटि 0।
\end{solution}

\begin{solution}{exo:b1:rate-laws:5}
$[\ce{A}]$, $t$ में रैखिक नहीं है (हानियाँ 0.0259, 0.0192, 0.0142, 0.0106)।
$\ln[\ce{A}] = -2.303$, $-2.602$, $-2.902$, $-3.202$, $-3.503$: यह हर
\qty{10}{min} में 0.300 घटता है, एक सरल रेखा। कोटि 1, $k =
\qty{0.0300}{min^{-1}}$।
\end{solution}

\begin{solution}{exo:b1:rate-laws:6}
$[\ce{A}]_0$ को दोगुना करने से $v_0$, $4 = 2^2$ से गुणा होता है: \ce{A} में कोटि 2।
$[\ce{B}]_0$ को दोगुना करने से वह 2 से गुणा होता है: \ce{B} में कोटि 1। $v =
k[\ce{A}]^2[\ce{B}]$, जहाँ $k = 2.0 \times 10^{-6}/(0.010^2 \times 0.010) =
\qty{2.0}{L^2.mol^{-2}.s^{-1}}$।
\end{solution}

\begin{solution}{exo:b1:rate-laws:7}
\ce{B} सौ गुना आधिक्य में है: उसकी सांद्रता
\qty{0.50}{mol/L} बनी रहती है, और $v = k'[\ce{A}]$, जहाँ $k' = k[\ce{B}]_0$। अतः
$k = 0.012/0.50 = \qty{0.024}{L.mol^{-1}.s^{-1}}$।
\end{solution}

\begin{solution}{exo:b1:rate-laws:8}
$E_a = R\ln5/(1/300 - 1/320) = 8.314 \times 1.609/(2.083 \times 10^{-4}) =
\qty{64.2}{kJ/mol}$। \qty{310}{K} पर: $\ln k = \ln(1.0 \times 10^{-3}) +
\frac{E_a}{R}\big(\frac{1}{300} - \frac{1}{310}\big) = -6.908 + 0.831$,
$k = \qty{2.3e-3}{s^{-1}}$।
\end{solution}

\begin{solution}{exo:b1:rate-laws:9}
यहाँ $a = 2$: $1/[\ce{A}] = 1/a_0 + 2kt$, $t_{1/2} = 1/(2k a_0) = 1/(2
\times 0.50 \times 0.020) = \qty{50}{s}$। \qty{100}{s} पर: $1/[\ce{A}] = 50
+ 100 = \qty{150}{L/mol}$, $[\ce{A}] = \qty{6.7e-3}{mol/L}$।
\end{solution}

\begin{solution}{exo:b1:rate-laws:10}
$[\ce{A}] = a_0(1 - f) = a_0\eu^{-akt_f}$ से $t_f = -\ln(1-f)/(ak)$ मिलता है, और
\omterm{def:b1:rate-laws:half-life}{अर्ध-आयुओं} में $t_f/t_{1/2} = -\ln(1-f)/\ln2$: 50\,\% के लिए 1, 75\,\% के लिए 2,
90\,\% के लिए 3.32, 99.9\,\% के लिए 9.97: अभिक्रिया के ``पूरी'' होने में लगभग दस
\omterm{def:b1:rate-laws:half-life}{अर्ध-आयु}।
\end{solution}

\begin{solution}{exo:b1:rate-laws:11}
$m(t) = 10 - 0.40\,t$ (\unit{mg}, \unit{h}): कोटि 0। \omterm{def:b1:rate-laws:half-life}{अर्ध-आयु}
$10/(2 \times 0.40) = \qty{12.5}{h}$; \qty{25}{h} के बाद ख़ाली। प्रति घंटे दी जाने वाली
मात्रा बची हुई मात्रा पर निर्भर नहीं करती: एक स्थिर आपूर्ति, जो
पैच का उद्देश्य ही है।
\end{solution}

\begin{solution}{exo:b1:rate-laws:12}
$E_a = R\ln2/(1/298.15 - 1/308.15) = 8.314 \times 0.693/(1.088 \times
10^{-4}) = \qty{52.9}{kJ/mol}$। 273.15 और \qty{283.15}{K} के बीच:
$\exp\big(\frac{52900}{8.314}(\frac{1}{273.15} - \frac{1}{283.15})\big) =
2.3$: वही \omterm{def:b1:rate-laws:activation-energy}{सक्रियण ऊर्जा} कम ताप पर बड़ा गुणक देती है।
\end{solution}

\begin{solution}{pb:b1:rate-laws:1}
\textbf{1.} हानियाँ 8.6, 7.9, 7.2 और 6.5 अंक: स्थिर नहीं, इसलिए
कोटि 0 नहीं।
\textbf{2.} $\ln$: 4.605, 4.515, 4.425, 4.335, 4.245; यह हर 10 दिन में 0.090
घटता है: एक सरल रेखा, कोटि 1।
\textbf{3.} $k_{60} = 0.090/10 = \qty{9.0e-3}{d^{-1}}$।
\textbf{4.} $t_{1/2} = \ln2/k_{60} = \qty{77}{d}$।
\textbf{5.} $100\,\eu^{-0.54} = 58.3\,\%$।
\textbf{6.} $t_{90} = \ln(10/9)/k_{60} = 0.1054/0.0090 = \qty{11.7}{d}$।
\textbf{7.} 90\,\% शेष: $\eu^{-kt_{90}} = 0.9$, इसलिए $t_{90} = \ln(10/9)/k$।
\textbf{8.} $t_{90}/t_{1/2} = \ln(10/9)/\ln2 = 0.152$।
\textbf{9.} कोटि 1 के लिए शेष भिन्न केवल $kt$ पर निर्भर करता है,
प्रारंभिक मात्रा पर नहीं।
\textbf{10.} \qty{25}{\celsius} पर निम्नीकरण में वर्षों लगते हैं; गर्म करना
इसे त्वरित कर देता है (आरेनियस), ताकि इसे सप्ताहों में मापा जा सके।
\textbf{11.} कोटि 2 के लिए \omterm{def:b1:rate-laws:half-life}{अर्ध-आयु} $1/(ak\,a_0)$ है: दोगुनी
प्रबल गोली आपेक्षिक रूप से दोगुनी तेज़ी से निम्नीकृत होती और कम समय टिकती।
\textbf{12.} $1/T$ (\unit{K^{-1}}): \num{3.1934e-3}, \num{3.0945e-3},
\num{3.0017e-3}; $\ln k$: $-6.669$, $-5.661$, $-4.711$।
\textbf{13.} प्रवणता $= (-4.711 + 6.669)/(3.0017 \times 10^{-3} - 3.1934 \times
10^{-3}) = \qty{-1.021e4}{K}$।
\textbf{14.} $E_a = 8.314 \times 1.021 \times 10^4 = \qty{85}{kJ/mol}$।
\textbf{15.} अनुमानित $\ln k_{50} = -4.711 - 1.021 \times 10^4 \times (3.0945 -
3.0017) \times 10^{-3} = -5.658$, मापा गया $-5.661$: रेखा पर।
\textbf{16.} $A = k_{60}\,\eu^{E_a/RT} = 0.0090\,\eu^{10215/333.15} =
\qty{1.9e11}{d^{-1}}$।
\textbf{17.} $\exp\big(1.021 \times 10^4(\frac{1}{298.15} -
\frac{1}{308.15})\big) = 3.0$: नियम के गुणक 2 से अधिक, क्योंकि
$E_a$, \qty{53}{kJ/mol} से बड़ी है।
\textbf{18.} $k_{30} = k_{60}\exp\big(-1.021 \times 10^4(\frac{1}{303.15} -
\frac{1}{333.15})\big) = \qty{4.3e-4}{d^{-1}}$।
\textbf{19.} $t_{90} = 0.1054/4.33 \times 10^{-4} = \qty{243}{d}$, 8.0
माह।
\textbf{20.} $1 - \eu^{-3 \times 0.0090} = 2.7\,\%$।
\textbf{21.} भंडारण-आयु ताप के साथ तेज़ी से घटती है (\qty{30}{\celsius} पर 8 माह,
\qty{25}{\celsius} पर 14 माह के मुक़ाबले): समाप्ति-तिथि
केवल \qty{25}{\celsius} से नीचे मान्य है।
\textbf{22.} $k_{25} = k_{60}\exp\big(-1.021 \times 10^4(\frac{1}{298.15} -
\frac{1}{333.15})\big) = \qty{2.46e-4}{d^{-1}}$।
\textbf{23.} $t_{90} = 0.1054/2.46 \times 10^{-4} = \qty{428}{d}$, अर्थात्
\textbf{14 माह}: डिब्बे पर छपी तिथि, जो छह
सप्ताह के मापनों से प्राप्त हुई।
\end{solution}
